%---------------------------------------------------------------------
\chapter{Conclusion and Future Work}
%---------------------------------------------------------------------

\section{Conclusion}

This report addressed the problem of checking whether a GI-tag claim on an Indian handloom product is
supported by official evidence, a task for which the necessary information is spread across the GI
registry, GI descriptions, the national handloom census and producer records. It presented LoomTrace,
an agent implemented as a six-node LangGraph state graph with seven typed tools, retrieval over official
GI text, a bounded re-planning loop, an evidence grader, a citation and confidence guardrail, and a
human review workflow, covering fifteen registered handloom GIs.

The objectives stated in Section~1.4 were addressed as follows.

\begin{enumerate}[itemsep=2pt]
\item \textbf{Agent design.} Met. The state-graph agent returns one of three verdicts with the result
and source of every check (Chapter~4).
\item \textbf{Tools, retrieval and bounded re-planning.} Met. Seven typed tools and a BM25 index over
21 passages were implemented; the loop is limited to three iterations and in practice used at most two
(Chapter~6).
\item \textbf{Guardrail and human review.} Met. All 15 unclear cases in the test set were escalated, and
reviewer decisions with notes are recorded and shown to the requester (Chapters~5 and~7).
\item \textbf{Tracing and checkpoints.} Met. Every node writes a trace entry and a checkpoint; LangSmith,
proposed in the design document, was replaced by this local store (Chapter~5).
\item \textbf{Evaluation against a baseline.} Met for accuracy and abstention. On the 45-case golden set
the agent reached 100\% accuracy and 100\% correct abstention against targets of 80\% and 90\%, while a
single-lookup baseline reached 33.3\% and 0\%. Citation faithfulness was not measured, and the result
holds for team-constructed cases only; on six realistic spelling variants the agent produced four wrong
\emph{Conflicting} verdicts (Chapter~7).
\item \textbf{A/B experiment.} Met. Disabling re-planning reduced accuracy from 100\% to 80\% and
increased escalations from 15 to 24, while the loop cost 0.8 tool calls and about 3\,ms per case;
re-planning was adopted as the default (Chapter~8).
\end{enumerate}

The main conclusion is that, for this verification task, the agentic elements contribute through
control rather than through generation: re-planning recovers cases that a single pass would wrongly
leave undecided, and the guardrail ensures that undecidable cases reach a person instead of receiving
a guessed verdict.

\section{Future Work}

The following extensions follow directly from the limitations in Chapter~9.

\begin{enumerate}[itemsep=2pt]
\item \textbf{Robust name matching.} Add fuzzy matching of product, district and state names with a
similarity band in which checks are marked \emph{unknown} rather than \emph{fail}, together with tables
of state abbreviations and former state names, so that the four failures in Table~\ref{tab:probes}
become correct or escalated verdicts.
\item \textbf{Official GI application text.} Download and OCR the statements of case of the covered GIs
from the IP India registry, index them with chunking and dense retrieval, and record the GI application
numbers of all fifteen products, replacing the team-written summaries and indicative page numbers.
\item \textbf{Real producer data.} Replace the simulated weaver and batch records with records from a
partner weavers' cooperative, with consent and access control, and evaluate on real labels.
\item \textbf{Independent evaluation.} Build a test set from real marketplace listings labelled by two
independent annotators, measure inter-annotator agreement and citation faithfulness, and evaluate the
language-model extraction path on it.
\item \textbf{Automatic resumption.} Use the stored checkpoints to resume interrupted runs and to pause
the graph at the review step, so that a reviewer's decision continues the same run.
\end{enumerate}
